% Build: cd docs/research/paper && tectonic cwmr-etf.tex   (writes cwmr-etf.pdf)
\documentclass[11pt]{article}

\usepackage[margin=1in]{geometry}
\usepackage[T1]{fontenc}
\usepackage{lmodern}
\usepackage{microtype}
\usepackage{amsmath,amssymb}
\usepackage{booktabs}
\usepackage[skip=6pt]{caption}
\usepackage[round]{natbib}
\usepackage[hidelinks]{hyperref}
\usepackage{xurl}


\title{Mean-Reversion Portfolio Selection on US-Listed ETFs, 1999--2026:\\
A Replication with Pre-Registered Tests}
\author{Yelnil Gabo}
\date{October 7, 2026\\[2pt] \small Working paper}

\begin{document}
\maketitle

\begin{abstract}
\citet{lahanis2025} report that online portfolio selection algorithms built on mean reversion,
CWMR and PAMR among them, earn large risk-adjusted returns on NASDAQ-100 stocks over 1998--2009.
We reproduce their tables exactly from their published code and data, then test the strongest
algorithm, the authors' variant of Confidence-Weighted Mean Reversion (CWMR), on 16 liquid
US-listed equity ETFs with realistic trading. Daily CWMR trades about 350 times its account a
year and trails SPY at measured 2016--2024 spreads. We searched variants that trade less,
selected one on 2016--2024 data and tested it on data it had not seen. A monthly variant
selected this way tied SPY on a pre-registered 2025--2026 holdout (18.28\% against 18.36\% a
year). Variants that learn from weekly or monthly bars passed a pre-registered 1999--2015 test
against SPY: weekly bars returned 9.82\% a year against SPY's 4.90\%. That test turned out to be
weak. An equal-weight portfolio of the same ETFs returned 6.61\%, and the strategy's lead over
it is not statistically significant ($p = 0.33$). The lead also disappears at trading costs
between 5 and 10 basis points per dollar traded. Before 2001 the minimum tick alone cost more
than that on sector ETFs, and we have no measured spreads for 2001--2010. We find no tradable
edge.
\end{abstract}

\section{Introduction}

Online portfolio selection treats investing as a repeated game. Each period an algorithm holds
a portfolio, observes the period's price changes and updates its weights by a fixed rule.
\citet{cover1991} showed that some rules provably approach the best constant-weight portfolio
in hindsight. A later family of algorithms bets on reversal instead. Anticor
\citep{borodin2004}, PAMR \citep{li2012pamr}, CWMR \citep{li2013cwmr} and RMR \citep{huang2013}
move weight from recent winners to recent losers. On the historical datasets in this literature
these algorithms report wealth multiples in the hundreds or thousands \citep{lihoi2014}.

\citet{lahanis2025} implement about 25 of these algorithms and test them on NASDAQ-100 stocks
from 1998 to 2009. Their best result, CWMR, multiplies wealth 1,590 times. They disclose that
they ignore trading costs and that the period may not represent recent markets.

This paper asks whether any version of their CWMR survives realistic trading on instruments a
retail investor can buy today. We use btest, a backtester that trades the way its live system
does, and we pre-register each confirmatory test before running it. Section~\ref{sec:data}
describes the data and methods. Section~\ref{sec:repro} reproduces the original paper.
Sections~\ref{sec:daily} to~\ref{sec:bars} report the out-of-sample tests in the order we ran
them, since the order matters for how much each result can be trusted.
Section~\ref{sec:deviations} reports where we departed from the pre-registered plans.
Section~\ref{sec:discussion}
discusses what the evidence supports.

\section{Data and methods}\label{sec:data}

\subsection{Algorithm}

All tests use the authors' CWMR as written in their public repository, ported
to btest and checked against their function to within $7.1\times10^{-9}$ per weight on their
data. Their CWMR differs from the published algorithm of \citet{li2013cwmr}: the step size has
no confidence term, a learning-rate factor $\eta$ is added and the update is not mean-centred.
We keep their parameters throughout ($\epsilon = 0.89$, $\theta = 0.92$, $\eta = 0.93$) and
never re-tune them.

Each period the algorithm receives a vector $x$ of price ratios, one per asset ($x_i = 1.02$
means asset $i$ rose 2\%). With current weights $b$ and covariance $\Sigma$ (initially the
identity), it computes
\begin{align}
  \lambda &= \eta \max\!\left(0,\ \frac{b^\top x - \epsilon}{x^\top \Sigma x}\right), \\
  \mu &= b - \lambda\, \Sigma x, \\
  \Sigma^{-1} &\leftarrow \Sigma^{-1} + 2 \lambda \theta\, x x^\top, \\
  b &\leftarrow \Pi_\Delta(\mu),
\end{align}
where $\Pi_\Delta$ projects onto non-negative weights that sum to one. Because $\epsilon = 0.89$,
the update fires unless the portfolio falls more than 11\% in one period, so in practice it
fires every period. Assets that rose most lose the most weight.

We test four schedules (Table~\ref{tab:schedules}).

\begin{table}[ht]
\centering
\caption{CWMR schedules tested.}\label{tab:schedules}
\begin{tabular}{lll}
\toprule
Name & Learns from & Trades \\
\midrule
Daily & daily ratios & every session \\
Monthly, daily signal & daily ratios & last session of each month \\
Monthly bars & month-end to month-end ratios & last session of each month \\
Weekly bars & week-end to week-end ratios & last session of each ISO week \\
\bottomrule
\end{tabular}
\end{table}

\subsection{Universe and data}

The universe is 16 ETFs: SPY, QQQ, IWM, DIA, EFA, EEM and ten Select Sector SPDRs (XLK, XLF,
XLV, XLE, XLI, XLY, XLP, XLU, XLB, XLRE). EFA and EEM hold non-US stocks and XLRE holds real
estate. The ETFs launched between 1993 and 2015, so the universe grows over time: it has 11
members in January 1999 and reaches 16 in October 2015. A fund joins on its first trading day.
Before that it holds no weight. Until a fund trades, the algorithm gives it the average price
ratio of the funds that do and splits its weight equally among them, which gives the same
portfolio return, so the algorithm runs on the whole universe from the start and keeps what it
has learned when a fund joins.

We use two price sources. From 2016 on, btest uses Alpaca SIP minute bars adjusted for splits
and dividends; a strategy decides at 15:30 New York time on data through 15:14 and fills at
15:45, which is how the live system trades. For runs that start before 2016 we use Yahoo daily
bars adjusted for dividends and splits; a strategy decides on a session's close and fills at
the next session's open. The stored Yahoo closes match a fresh download to within a ratio of
$3\times10^{-6}$ on every date.

\subsection{Costs and benchmarks}

Every dollar traded pays a flat cost in basis points (bp), plus the SEC fee on sales. We
measured half-spreads at 15:45 from Alpaca quotes on 45 sessions in 2016--2024
(Table~\ref{tab:spreads}).

\begin{table}[ht]
\centering
\caption{Median half-spreads at 15:45, 2016--2024.}\label{tab:spreads}
\begin{tabular}{lr}
\toprule
ETF & Half-spread (bp) \\
\midrule
SPY & 0.17 \\
QQQ, IWM, DIA & 0.26--0.30 \\
Most sector SPDRs, EFA & 0.41--0.82 \\
EEM, XLRE, XLF & 1.20--1.68 \\
Equal-weight average & 0.70 \\
\bottomrule
\end{tabular}
\end{table}

We charge 0.7~bp unless stated and report results at 5, 10 and 20~bp. The model has no
commissions, no market impact and no taxes; Section~\ref{sec:discussion} discusses each.

The benchmarks are SPY with dividends reinvested and an equal-weight portfolio of the same
ETFs, reset at each month end. The second matters more. A strategy that beats SPY only because
these ETFs beat SPY has learned nothing.

\subsection{Statistics and procedure}

We compare Sharpe ratios with the test of \citet{jobson1981} with the correction of
\citet{memmel2003}, on daily returns of two strategies over the same days. The test uses their
correlation, 0.71 to 0.95 across this study, so it can detect smaller differences than
comparing two independent estimates. It assumes independent returns. The daily return
differences have negative first-order autocorrelation, which makes the test conservative.

Each confirmatory test was written into the research log and committed to git before it ran.
The commit names the strategy version, window, settings and pass criterion, and states that the
strategy will not be changed and rerun on that window. Each commit precedes its run; the
holdout's commit, for example, is 28 seconds older than the run's database record. We kept
2025-01-01 onward as a holdout and have used it once.

\section{Reproducing the original paper}\label{sec:repro}

The authors' repository contains their data, 93 NASDAQ-100 stocks from 1998-01-02 to
2009-12-31. Running their functions reproduces their Tables~3 and~4 exactly: CWMR 1,590.38$\times$,
PAMR 788.26$\times$, Anticor 36.89$\times$, FTRL 15.25$\times$ and the equal-weight constant
rebalanced portfolio (CRP) 12.16$\times$.

The paper's parameters were chosen on this same sample, so these are in-sample results. Within
the sample, mean reversion keeps much of its return under frictions (Table~\ref{tab:paper}).

\begin{table}[ht]
\centering
\caption{The authors' data, 1998--2009. Paired Sharpe test against CRP.}\label{tab:paper}
\begin{tabular}{lrr}
\toprule
 & CAGR & $p$ vs CRP \\
\midrule
CRP & 23.2\% & \\
CWMR & 85.0\% & $<0.0001$ \\
CWMR, 10 bp & 43.6\% & 0.29 \\
CWMR, one-day delay & 45.6\% & 0.11 \\
PAMR, one-day delay & 51.9\% & 0.0002 \\
PAMR, one-day delay and 5 bp & 33.9\% & 0.49 \\
\bottomrule
\end{tabular}
\end{table}

Five of the paper's strategies differ from the algorithms they are named after. As implemented,
OLMAR and exponential gradient never change their target weights. Both reduce to CRP, and the
paper reports identical numbers for all three. FTRL's lead over CRP disappears when the 32
tickers with incomplete data are removed. CWMR and RMR are modified versions of the published
methods. Two bad prints, PCAR in February 2000 and XRAY in July 2004, account for about 30\% of
CWMR's final wealth.

\section{Daily trading, 2016--2024}\label{sec:daily}

On the 16 ETFs with btest's live timing, daily CWMR returns 16.6\% a year before costs and
13.8\% at 0.7~bp, against SPY's 14.5\%. It turns over about 350 times its account a year, so
each basis point of cost removes about 3.5 percentage points of annual return. Its break-even
cost against SPY lies between 0.43 and 0.70~bp. Paired tests against SPY are not significant at
any setting we tested ($p = 0.51$ for the paper's close-to-close timing before costs). Daily CWMR
is not worth pursuing at retail costs, and the remaining sections test schedules that trade less.

\section{Selecting a lower-turnover variant, and the holdout}\label{sec:holdout}

We ran all 13 of the paper's strategy families three ways on 2016--2024: daily, daily signal
with monthly trades, and monthly bars. CWMR with a daily signal and monthly trades did best,
18.3\% a year (Sharpe 0.87) against SPY's 14.5\% (0.73), with turnover cut to 18$\times$. Its
paired $p$ against SPY was 0.24, and it was the best of about 26 variants, so we treated it as a
hypothesis and pre-registered a test on the untouched 2025--2026 holdout. It tied SPY
(Table~\ref{tab:holdout}).

\begin{table}[ht]
\centering
\caption{Pre-registered holdout, 2025-01-01 to 2026-10-02.}\label{tab:holdout}
\begin{tabular}{lrr}
\toprule
 & Monthly CWMR & SPY \\
\midrule
CAGR & 18.28\% & 18.36\% \\
Sharpe (T-bill) & 0.81 & 0.87 \\
Max drawdown & $-20.3$\% & $-18.7$\% \\
\bottomrule
\end{tabular}
\end{table}

The paired $p$ was 0.77, and the strategy beat SPY in 8 of 21 months. A later continuous run
showed why the 2016--2024 result was misleading. The 18.3\% depends on starting the algorithm
fresh in January 2016. Run on daily ETF data from 1995 with its state carried forward, the same
strategy returns 11.2\% over 2016--2024 against SPY's 14.4\%.

\section{Learning from weekly and monthly bars, 1999--2024}\label{sec:bars}

The monthly variant above learns from daily moves but holds for a month. Short-horizon reversal
in stock returns is documented over weeks \citep{lehmann1990} and a month \citep{jegadeesh1990},
so a signal built from one day's moves and held for a month measures one horizon and bets on
another. Variants that learn from week-end or month-end prices remove that mismatch. Over
2016--2024 they returned 17.6\% (weekly bars) and 11.6\% (monthly bars) against SPY's 14.5\%.
Neither had been run before 1999.

\subsection{Pre-registered test, 1999--2015}

We pre-registered both variants on 1999--2015 with ETF-only daily data. The pass criterion was
CAGR and Sharpe ratio above SPY's at 0.7~bp. This period contains two crashes, and SPY returned
only 4.9\% a year (Table~\ref{tab:prereg}).

\begin{table}[ht]
\centering
\caption{Pre-registered test, 1999--2015, 0.7 bp. Paired Sharpe tests.}\label{tab:prereg}
\begin{tabular}{lrrrrrr}
\toprule
 & CAGR & Sharpe & Max DD & Turnover & $p$ vs SPY & $p$ vs EW \\
\midrule
Weekly bars & 9.82\% & 0.45 & $-50.1$\% & 79$\times$ & 0.11 & 0.33 \\
Monthly bars & 8.90\% & 0.40 & $-61.4$\% & 14$\times$ & 0.17 & 0.52 \\
Equal weight (EW) & 6.61\% & 0.33 & $-54.4$\% & 0.3$\times$ & 0.10 & \\
SPY & 4.90\% & 0.25 & $-55.2$\% & 0$\times$ & & \\
\bottomrule
\end{tabular}
\end{table}

Both variants pass the criterion as written. But equal weight also passes it, and the paired
tests cannot separate either variant from equal weight. In hindsight the criterion was too
easy. We already knew this basket of ETFs had beaten SPY over this window when we wrote it, so
equal weight was the right benchmark to name.

\subsection{Robustness to the day the week ends}

Weekly bars end on Friday. We reran the strategy with 5-session blocks ending on each of the
other four weekdays, each version once over 1999--2024 (Table~\ref{tab:offsets}).

\begin{table}[ht]
\centering
\caption{Weekly bars by week-end day, 0.7 bp. CAGR.}\label{tab:offsets}
\begin{tabular}{lrrr}
\toprule
 & 1999--2015 & 2016--2024 & 1999--2024 \\
\midrule
Calendar weeks (Friday) & 9.9\% & 16.4\% & 12.1\% \\
Five shifted versions & 7.8--13.4\% & 11.6--14.8\% & 9.6--13.5\% \\
Equal weight & 6.6\% & 11.8\% & 8.4\% \\
SPY & 4.9\% & 14.4\% & 8.1\% \\
\bottomrule
\end{tabular}
\end{table}

Before 2016 every version beats both benchmarks. After 2016 only the Friday version and one
shifted version beat SPY. Over the whole window the Friday version is the second best of six,
and its 2016--2024 result is the best, which is what selection on 2016--2024 would produce.

The weekly variant is also concentrated. After rebalancing it holds a median of 2 ETFs, and one
ETF carries 90\% or more of the account in 36\% of weeks.

\subsection{Costs}

At about 76$\times$ turnover, each basis point of cost removes about 0.8 percentage points a
year (Table~\ref{tab:costs}).

\begin{table}[ht]
\centering
\caption{Weekly bars at higher costs per dollar traded. CAGR.}\label{tab:costs}
\begin{tabular}{lrrrrr}
\toprule
 & 0.7 bp & 5 bp & 10 bp & 20 bp & SPY \\
\midrule
1999--2015 & 9.82\% & 6.24\% & 2.23\% & $-5.35$\% & 4.90\% \\
1999--2024 & 12.10\% & & 4.46\% & & 8.09\% \\
\bottomrule
\end{tabular}
\end{table}

The break-even against SPY is about 6--7~bp over 1999--2015 and about 5--6~bp over 1999--2024.
Before April 2001 US stocks traded in sixteenths of a dollar, so the smallest possible
half-spread on a \$30 sector SPDR was about 10~bp. We found no measured spreads for these ETFs
in 1999--2005 and cannot say when they fell below 5~bp. Over 1999--2001 the strategy also made
243 to 315 trades a year. At the 1999 average online commission of about \$16 a trade, that is
about a fifth of a \$20,000 account each year.

\section{Deviations from the pre-registered plans}\label{sec:deviations}

Two results differ from what was pre-registered. First, the 1999--2015 test of
Section~\ref{sec:bars} was first run with each order filling at the closing price its decision
had used, which no one can trade at. We moved fills to the next session's open and reran both
variants once with nothing else changed. Before the correction weekly bars returned 11.50\% a
year (paired $p = 0.02$ against SPY) and monthly bars 8.47\%. Both passed the criterion before
and after; Table~\ref{tab:prereg} reports the corrected run.

Second, an earlier pre-registered test of the monthly variant on 1995--2015 filled the years
before each ETF launched with similar mutual funds and indexes, and used the same uncorrected
fills. It passed, returning 11.57\% a year against SPY's 9.33\%. We no longer rely on it. With
each ETF's own history only and corrected fills, the same strategy returns 6.79\% over that
window. It holds only SPY until DIA launches in 1998 and sits in cash before then.

\section{Discussion}\label{sec:discussion}

The original paper's results are real, in the sense that its code computes them correctly from
its data. They do not survive the move to instruments and costs a retail investor faces today.
Every schedule we tried either trails SPY at measured costs (daily), ties it out of sample
(monthly, on the holdout), or beats SPY by roughly as much as an equal-weight portfolio of the
same ETFs does (weekly and monthly bars, 1999--2015).

The weekly-bars variant is the most interesting result. Before 2016 it beat both benchmarks at
every week-end offset. Several explanations fit. Weekly reversal across sector ETFs may have
been real in a market with wider spreads and slower arbitrage, in which case the profit belonged
to whoever supplied liquidity and was not available to someone paying the spread. Daily returns
of these ETFs in excess of SPY show negative next-day autocorrelation before 2016 (median across
the ETFs $-0.065$ in 1999--2003, $-0.092$ in 2004--2008 and $-0.047$ in 2009--2015) and none
after (0.000 in 2016--2024). That pattern fits bid-ask bounce in closing prints,
which a backtest captures and a trader pays for. And the variant came out of a search over
about 100 configurations across this study, which inflates the best result
\citep{harvey2016}. We cannot separate these with the data we have.
Each predicts the effect fades as markets become cheaper to trade, and after 2016 it roughly ties
SPY.

Before calling it an edge we would need:
\begin{enumerate}
  \item measured spreads for these ETFs in 1999--2010, to replace the flat cost;
  \item a pre-registered test against equal weight, not SPY, on data no variant has touched. The
        2025 holdout has been used once; data from 2027 on would be clean;
  \item results after commissions for a stated account size, and after tax for taxable
        accounts. At 76$\times$ turnover almost every gain is short-term.
\end{enumerate}

\section{Conclusion}

We reproduced a published set of mean-reversion results exactly, then lost them to costs, to
timing and to the benchmark. A weekly variant passed a pre-registered test against SPY, but the
test was weak: an equal-weight portfolio of the same ETFs passes it too, and the variant's edge
disappears at costs these ETFs carried twenty years ago. We find no evidence of a tradable edge
for this family of strategies on liquid US-listed ETFs.

\section*{Data and code availability}

Code that reproduces Section~\ref{sec:repro} and Tables~\ref{tab:paper}, \ref{tab:prereg},
\ref{tab:offsets} and~\ref{tab:costs} from public data is at
\begin{center}\url{https://github.com/yelgabo/cwmr-etf-replication}\end{center}
It uses the original authors' public code
and data \citep{lahanis2025}, Yahoo Finance daily prices and the 3-month Treasury bill rate from
FRED. Sections~\ref{sec:daily} and~\ref{sec:holdout} used minute bars from Alpaca, which require
an account; that code is available from the author on request.

\section*{Acknowledgments}

The author used Claude (Anthropic) to write analysis code and edit the manuscript.

\bibliographystyle{plainnat}
\bibliography{refs}

\end{document}
